Write $p=p_c$, $A=2/p_c$, $S(t)=e^{\mathrm it\Delta}$, and $I_0=[1/2,1]$. Fix $\chi\in C_c^\infty(B(0,2))$ equal to one on $B(0,1)$, and write $P_N=\chi(D/N)$, so that $P_Nf=f$. Tile $\mathbb R^d$ by cubes $Q$ of side $N$ and let $30Q$ be the concentric cube of side $30N$.

For $|z|\geq8N$ and $t\in I_0$, the kernel $K_{N,t}$ of $S(t)P_N$ satisfies
\[
|K_{N,t}(z)|\leq C_M N^d(1+N|z|)^{-M}.
\]
Indeed, after $\xi=N\eta$ the phase gradient in $\eta$ is $Nz-2tN^2\eta$, of magnitude at least $N|z|/2$ on the cutoff support. Its derivatives of positive order are either $O(N^2)$ or zero, and $N^2\lesssim N|z|$. Iterated integration by parts gives the uniform bound. For $x\in Q$ and $y\notin30Q$ we have $|x-y|\geq8N$, so the far contribution on every cube is bounded by the same convolution of $|f|$ with $C_M N^d(1+N|\cdot|)^{-M}$. For $M>d$ the latter kernel has bounded $L^1$ norm, giving a bound $O(\|f\|_p)$ on $\mathbb R^d\times I_0$ by Young's inequality.

For the near contribution translate $Q$ to the origin and put $h=f\mathbf 1_{30Q}$ in the translated coordinates. The two multipliers commute on $L^2$, and
\[
S(t)P_Nh(x)=P_NS(t)h(x)
 =\int\check\chi(w)\,S(t)h(x-w/N)\,\mathrm dw.
\]
Here $\check\chi$ is the inverse Fourier transform with the factor $(2\pi)^{-d}$. Take $\lambda=16\sqrt d\,N$ and $g(v)=h(\lambda v)$, which is supported in $B_d=B(0,1)\subset\mathbb R^d$. Since $h\in L^1\cap L^2$, the Fresnel formula applies and its spatial integral is absolutely convergent. Making the substitution $y=\lambda v$ yields
\[
|S(t)h(z)|=(4\pi t)^{-d/2}\lambda^d
 \left|Eg\!\left(-\frac{\lambda z}{2t},\frac{\lambda^2}{4t}\right)\right|,
\quad
Eg(x,t)=\int_{B_d}e^{\mathrm i(x\cdot v+t|v|^2)}g(v)\,\mathrm dv.
\]
For $z=x-w/N$ with $x\in Q$ and $t\in I_0$, the two arguments lie together in a ball of radius $C_d(N^2+|w|)$. The map $(x,t)\mapsto(-\lambda(x-w/N)/(2t),\lambda^2/(4t))$ is one-to-one by its last coordinate. Its absolute Jacobian is $(\lambda/(2t))^d\lambda^2/(4t^2)\asymp_d\lambda^{d+2}$. Since $\|g\|_p=\lambda^{-d/p}\|h\|_p$, \ref{ms:lemma-12-cap} gives
\[
\|S(t)h(x-w/N)\|_{L^p(Q\times I_0)}
 \lesssim_d \lambda^{d-(2d+2)/p}[\log(2+N^2+|w|)]^A\|h\|_p.
\]
The power of $\lambda$ is zero. By Minkowski's inequality and the rapid decay of $\check\chi$, the near contribution has norm at most $C[\log(2+N)]^A\|f\|_{L^p(30Q)}$. Sum its $p$th powers over $Q$ using the bounded overlap of $30Q$, and add the far bound. This proves the claim on $I_0$.

On $[a/2,a]$, $a>0$, set $f_a(x)=f(\sqrt a\,x)$. Then $S(t)f_a(x)=(S(at)f)(\sqrt a\,x)$ and the support radius for $\widehat f_a$ is at most $\sqrt a\,N$. Scaling the estimate on $I_0$, the spatial factors cancel and leave
\[
\|S(t)f\|_{L^p(\mathbb R^d\times[a/2,a])}
 \lesssim_d a^{1/p}[\log(2+\max(2,\sqrt a\,N))]^A\|f\|_p.
\]
For bounded positive times sum the $p$th powers over intervals with $a=T2^{-j}$, $j\geq0$, and fixed $T>0$. The logarithm is bounded by $C_T\log(2+N)$, and $\sum_j T2^{-j}<\infty$. For negative times use $S(-t)f=\overline{S(t)\overline f}$, with the same frequency-ball condition. This gives the result on $I$.
